\input amstex
\documentstyle{amsppt}
\magnification 1200
\hsize 6.5 true in
\vsize 8.7 true in
\define\off{\overline{F}}
\define\ozf{\overline{Z(F)^n}}
\define\oglf{\overline{GL(2,F)}}
\define\oa{\overline{A(F)}}
\define\oan{\overline{A(F)_n}}
\define\oad{\overline{A(F)_d}}
\define\oslf{\overline{G(F)}}
\define\eps{\epsilon}
\define\okk{{\Cal O}_K}
\baselineskip 16pt
\subheading{Notes on the inflation-restriction sequence
in non-abelian cohomology}
\vskip .1in
Let $E/F$ denote a finite Galois extension of p-adic fields,
$G$ a connected reductive group over $F$, $\hat{G}$ its Langlands
dual over ${\Bbb C}$, and $\Gamma =Gal(E/F)$. Of course, if $G$ is
$E$-split but not $F$-split then $\Gamma$ acts on $\hat{G}({\Bbb C})$ 
non-trivially. The relative Weil group $W=W_{E/F}=W_F/W_E^c$ fits into the
short exact sequence
$$
1\rightarrow E^{\times}
\overset i\to \rightarrow W \overset p\to \rightarrow 
\Gamma\rightarrow 1,
$$
where $W_E^c$ denotes the closure of the commutator 
subgroup of the absolute Weil group $W_E$.
Note that $E^\times$ is isomorphic to $W_F/W_F^c$,
the maximal abelian quotient of $W_F$. This sequence allows us to 
pull-pack the $\Gamma$-action on $\hat{G}({\Bbb C})$ to
a $W$-action. The pull-back of this $W$-action to an action
of $E^\times$ is trivial since the sequence is (semi-)exact
at $W$.
\par
The purpose of this note is to verify the following
non-abelian analog of a sequence in \cite{Se, Ch VII, \S 6}. 
The proof is an exercise in cohomology but it is hoped that
the sequence itself may provide some insight into the
nature of Langlands parameters, which form a subset
of $H^1(W,\hat{G}({\Bbb C}))$. A good reference for this type of 
material is Springer's 1965 Boulder conference article \cite{Spr}.
\par
I thank A. DuPre for conversations on $H^2$.
\proclaim{Lemma} Let $s:\Gamma\rightarrow W$ be a section
for which $s(1)=1$ and $s(w^{-1}gw)=w^{-1}s(g)w$, for all
$g\in \Gamma$ and $w\in W$. Let
$$
\sigma(g_1,g_2)=s(g_1g_2)s(g_2)^{-1}s(g_1)^{-1},
$$
for $g_1,g_2\in \Gamma$.
\newline
(1) $\sigma :\Gamma\times\Gamma\rightarrow E^\times$
is a 2-cocycle representing a non-trivial class
in $H^2(\Gamma,E^\times)$, where $\Gamma$ acts on $E^\times$
by conjugation in $W$. Moreover,
$$
\sigma(g_1g_2,g_3)\sigma(g_1,g_2)=
\sigma(g_1,g_2g_3)\sigma(g_2,g_3)^{s(g_1)},
$$
for all $g_i\in \Gamma$. Here,
$$
\sigma(g_2,g_3)^{s(g_1)}=s(g_1)\sigma(g_2,g_3)s(g_1)^{-1}.
$$
\newline
(2) $(\partial f)(g_1,g_2)=f(\sigma(g_1,g_2))$, $g_1,g_2\in 
\Gamma$, defines a coboundary map making the following
sequence exact at every term but 
$H^1(E^\times,\hat{G}({\Bbb C}))^\Gamma$.
The sequence is semi-exact at 
$H^1(E^\times,\hat{G}({\Bbb C}))^\Gamma$.
$$
1\rightarrow H^1(\Gamma,\hat{G}({\Bbb C}))
\overset \text{inf}\to \rightarrow
H^1(W,\hat{G}({\Bbb C}))
\overset \text{res}\to \rightarrow
H^1(E^\times,\hat{G}({\Bbb C}))^\Gamma
\overset \partial\to \rightarrow
H^2(\Gamma,\hat{G}({\Bbb C}),\overline{\alpha}),
$$
where $H^2(\Gamma,\hat{G}({\Bbb C}),\overline{\alpha})$ is 
defined in the proof below, $\alpha$ denotes the action of 
$\Gamma$ on $\hat{G}({\Bbb C})$ and $\overline{\alpha}$ denotes
its class in the outer automorphism group (in the category
of algebraic ${\Bbb C}$-groups) of $\hat{G}({\Bbb C})$
$Aut(\hat{G}({\Bbb C}))/\hat{G}({\Bbb C})$. 
\endproclaim
\demo{Remark} This sequence in (2) holds for any algebraic
${\Bbb C}$-group $A$ with an algebraic 
$\Gamma$-action. In fact, we replace $\hat{G}({\Bbb C})$ by 
an arbitrary algebraic $\Gamma$-group $A$ defined over 
${\Bbb C}$ from now on. 
\enddemo
We now prove the lemma.
\par
(1) Since the section $s$ picks a representative in 
$W$ from each coset $wE^\times \in W/E^\times \cong\Gamma$,
$\in E^\times$ and so $s(g_1g_2)s(g_2)^{-1}s(g_1)^{-1}\in 
E^\times$. We write out $s(g_1g_2g_3)$ in two ways using
associativity and compare expressions:
$$
\split
s(g_1(g_2g_3))
&=\sigma(g_1,g_2g_3)s(g_1)s(g_2g_3)\\
&=\sigma(g_1,g_2g_3)s(g_1)\sigma(g_2,g_3)s(g_2)s(g_3)\\
&=\sigma(g_1,g_2g_3)\sigma(g_2,g_3)^{s(g_1)}s(g_1)s(g_2)s(g_3).
\endsplit
$$
On the other hand,
$$
\split
s((g_1g_2)g_3)
&=\sigma(g_1g_2,g_3)s(g_1g_2)s(g_3)\\
&=\sigma(g_1g_2,g_3)\sigma(g_1,g_2)s(g_1)s(g_2)s(g_3).
\endsplit
$$
Comparing both sides gives the identity.
\par
(2) We define terms an notions as they are needed.
\vskip .1in
{\bf $H^0$ and $H^1$}
\vskip .1in
Recall
$$
H^0(\Gamma,A)=A^\Gamma=\{a\in A\ |\ ga=a,\forall g\in\Gamma\}.
$$
A {\bf 1-cocycle} is a map $f:\Gamma\rightarrow A$ such that
$$
f(g_1g_2)=f(g_1)g_1f(g_2),
$$
for all $g_1,g_2\in \Gamma$. The set of such cocyles is denoted
$Z^1(\Gamma,A)$. Two 1-cocycles are called {\bf cohomologous}
if there is an $a\in A$ such that 
$$
f(g)=a^{-1}f'(g)g(a),
$$
for all $g\in \Gamma$. We write $f\sim f'$ if $f$ is cohomologous 
to $f'$. This is an equivalence relation.          
The first cohomology set of $\Gamma$ with values 
(or ``coefficients'') in $A$ is the set $H^1(\Gamma,A)$
of equivalence classes of elements in $Z^1(\Gamma,A)$.
The equivalence class of the constant cocycle $f=1$ is 
a distinguished point in this set.
\vskip .1in
{\bf Inflation} 
\vskip .1in
The homomorphism $r:W\rightarrow \Gamma$
defines a morphism of the functor
$$
A\longmapsto H^0(\Gamma,A)=A^\Gamma,
$$
on the category of $A$-groups, to the functor
$$
A\longmapsto H^0(W,p^*A)=(p^*A)^W,
$$
on the category of $W$-groups, where $p^*A$ denotes the group $A$ 
with the $W$-action induced by $p$. This ``extends'' to a 
morphism of functors, called inflation,
$$
p_1^*:H^1(\Gamma,A)\rightarrow H^1(W,p^*A),
$$
by pulling back cocycles via $p$: $p_1^*:f(g) \longmapsto
p_1^*f(w)=f(p(w))$. We have 
$$
p_1^*f(w_1w_2)=f(p(w_1)p(w_2))=
f(p(w_1))p(w_1)f(p(w_2))
=p_1^*f(w_1)w_1(p_1^*f)(w_2),
$$
where $w_i\in W$ and the action of $W$ on $p^*A$ is via $p$.
Moreover, $p_1^*$ preserves equivalence classes: if there exists
an $a\in A$ such that $f(g)=a^{-1}f'(g)ga$, for all $g\in 
\Gamma$, then $p_1^*f(w)=a^{-1}p_1^*f'(w)wa$, where $p(w)=g$
and $wa=p(w)a=ga$.
\par
Next we show that $p_1^*$ is injective, as a map of pointed sets.
Suppose $p_1^*f,p_1^*f'\in Z^1(W,p^*A)$ are cohomologous,
so there is an $a\in p^*A$ for which
$p_1^*f(w)=a^{-1}p_1^*f'(w)wa$, for all
$w\in W$. If $g=p(w)$ then $f(g)=a^{-1}f'(g)ga$.
Since $p$ is surjective, this forces $f$ and $f'$ to be
cohomologous. $p_1^*$ clearly sends the distinguished point,
the trivial class, of $H^1(\Gamma,A)$ to that of
$H^1(W,p^*A)$. Therefore, $p_1^*$ is injective.
\vskip .1in
{\bf Restriction}
\vskip .1in
Since
$$
H^0(W,p^*A)=(p^*A)^W,\ \ \ \ \ \ \
H^0(E^\times,i^*A)=i^*A,
$$
for any $\Gamma$-group $A$, the injection 
$i:E^\times \rightarrow W$ induces a morphism of functors
$$
i_0^*:H^0(W,p^*A)\rightarrow H^0(E^\times,i^*A)^\Gamma.
$$
We extend this restriction functor to a morphism
of functors
$$
i_1^*:H^1(W,p^*A)\rightarrow H^1(E^\times,i^*A)^\Gamma,
$$
as follows. Let $i^*f(x)=f(i(x))$, for $x\in E^\times$,
$f\in Z^1(W,p^*A)$.
Since $E^\times$ acts on $i^*A$ trivially and since
$f(w_1w_2)=f(w_1)w_1f(w_2)$, $w_i\in W$,
we have $i^*f(x_1x_2)=
f(i(x_1)i(x_2))=i^*f(x_1)x_1(i^*f)(x_2)
=i^*f(x_1)i^*f(x_2)$, for all $x_j\in E^\times$, $j=1,2$.
This proves $i^*f\in Z^1(E^\times, i^*A)$, if
$f\in Z^1(W, p^*A)$. To see that $i^*$ preserves equivalence 
classes, suppose $f,f'\in Z^1(W,p^*A)$ are cohomologous: there
is an $a\in A$ such that
$$
f(w)=a^{-1}f'(w)wa,\ \ \ \ \ w\in W.
$$
Then $i^*f(x)=a^{-1}i^*f'(x)a$, if $x\in E^\times$. This implies 
that $i^*f$ is cohomologous to $i^*f'$. Therefore, the 
restriction map $res=i_1^*$ is a well-defined map on
cohomology classes,
$$
res:H^1(W,p^*A)\rightarrow H^1(E^\times, i^*A).
$$
\proclaim{Claim} The action of $W$ on functions 
$f:E^\times\rightarrow A$,
$$
w^*:f(x)\longmapsto w^*f(x)=f(wx),
$$
defines a map $w^*:Z^1(E^\times,i^*A)\rightarrow 
Z^1(E^\times,i^*A)$ which preserves equivalence classes.
Moreover, $w^*$ only depends on its coset $wE^\times$.
Let $H^1(E^\times,i^*A)^W=H^1(E^\times,i^*A)^\Gamma$ denote the classes fixed
by every such $w^*$, $w\in W$. Then the image of the
restriction map lies in the $\Gamma$-invariant cocycles:
$$
res:H^1(W,p^*A)\rightarrow H^1(E^\times, i^*A)^\Gamma
=Hom_\Gamma(E^\times,i^*A)/A,
$$
where the quotient is via the action of $A$
on $Hom_\Gamma(E^\times,i^*A)$ given by
$a:f(x)\longmapsto a^{-1}f(x)a$.
\endproclaim
\demo{proof of claim} If $f\in Z^1(W,p^*A)$, $w\in W$ and $x\in 
E^\times$, then $w^*i^*f(x)=f(wi(x))=f(w)\cdot wi^*f(x)
=w^*f(1)\cdot wi^*f(x)$. Now, one can easily show that the 
cocycle condition for $f$ implies 
$w^*f(x_1x_2)=w^*f(x_1)w^*f(x_2)$, for all $x_i\in E^\times$.
Taking $x_1=1$, we find that $w^*f(1)=1$. This implies that
$w^*i^*f(x)=wi^*f(x)$, for $x\in E^\times$.
Since $H^1(E^\times, i^*A)=Hom(E^\times,i^*A)/A$, the only
thing left to do is to recall that $Hom_\Gamma(E^\times,i^*A)=
Hom(E^\times,i^*A)^\Gamma$, by definition. \qed (claim)
\enddemo
\vskip .1in
{\bf Semi-exactness at $H^1(W,p^*A)$}
\vskip .1in
We want to show that, for any $f\in H^1(\Gamma,A)$,
$res(inf(f))\in  H^1(E^\times, i^*A)$ is the trivial class.
This follows from the semi-exactness of the short exact sequence
for $W$:
$$
res(inf(f))=f(p(i(x)))=f(1),
$$
for all $x\in E^\times$. Since $H^1(E^\times,i^*A)=
Hom(E^\times,i^*A)/A$, $res(inf(f))(x)$ is a homomorphism.
This forces $f(1)=1$, which forces $res(inf(f))$ to 
represent the trivial class.
\vskip .1in
{\bf Exactness at $H^1(W,p^*A)$}
\vskip .1in
Let $f\in H^1(W,p^*A)$ belong to the kernel of $res$,
so $res(f)\in H^1(E^\times,i^*A)^\Gamma$
\newline
$ = Hom_\Gamma(E^\times,i^*A)/A$ is the trivial class. We want to show 
that there is an $f'\in H^1(\Gamma,A)$ such that $f=inf(f')$.
By assumption, our class $f$ is represented by a 1-cocycle,
also denoted $f$, satisfying $f(x)=1$ for all $x\in E^\times$.
The cocycle condition and $f(x)=1$ on $E^\times$
imply that $f(wx)=f(w)$, for all $w\in W$ and $x\in E^\times$.
In other words, $f$ factors through to a function on
the quotient $W/E^\times\cong \Gamma$, also denoted $f$.
To check that $f:\Gamma\rightarrow A$ satisfies the 
cocycle condition, let $g_1,g_2\in \Gamma$ and
$w_1,w_2\in W$ be related by $p(w_j)=g_j$, $j=1,2$.
Then
$$
f(g_1g_2)=f(w_1w_2E^\times)=f(w_1)w_1f(w_2)
=f(g_1)g_1f(g_2),
$$
by the cocycle condition on $W$. This verifies exactness.
\vskip .1in
{\bf The connecting map $\partial$}
\vskip .1in
Given $f\in Z^1(E^\times,i^*A)^\Gamma$, we define
$\partial f$ by
$$
\partial f(g_1,g_2)=f(\sigma(g_1,g_2)),
$$
for $g_1,g_2\in \Gamma$, where $\sigma$ is given as in part (1)
of the lemma. Following Springer, a 2-cocycle of $\Gamma$ with 
values in $A$ to be a triple $(f,\alpha,\overline{\alpha})$,
where $f:\Gamma\times\Gamma\rightarrow A$ is a map for 
which
$$
f(g_1g_2,g_3)f(g_1,g_2)=f(g_1,g_2g_3)\alpha(g_1)f(g_2,g_3),
$$
for all $g_i\in\Gamma$, 
$\overline{\alpha}:\Gamma\rightarrow Aut(A)/A$ is a fixed 
outer automorphism ($A$ acts on $Aut(A)$ by conjugation), 
and ${\alpha}:\Gamma\rightarrow Aut(A)$ is a fixed 
$\Gamma$-group structure on $A$ for which the class of 
$\alpha(g)$ in $Aut(A)/A$ is $\overline{\alpha}(g)$, for all 
$g\in G$. The set of all such 2-cocycles is denoted
$Z^2(\Gamma,A,\overline{\alpha})$.
We must show that $\partial f$ satisfies 
this identity, where $\alpha$ denotes the given action
of $\Gamma$ on $A$. 
Observe that if $p(w)=g$, for $w\in W$ and $g\in 
\Gamma$, and if $x\in E^\times$ then
$w^{-1}i(x)w=i(gx)$ (where of course $gx$ represents the Galois action
of $g$ on $x$). Therefore, applying $f\in Hom_\Gamma(E^\times,i^*A)$
to both sides of the 2-cocycle identity in part (1),
we find that
$$
\partial f(g_1g_2,g_3)\partial f(g_1,g_2)=\partial f(g_1,g_2g_3)
g_1\partial f(g_2,g_3).
$$
The last step used the fact that
$f(\sigma(g_2,g_3)^{s(g_1)})=
f(s(g_1)\sigma(g_2,g_3)s(g_1)^{-1})$
\newline
$=f(g_1\sigma(g_2,g_3))=g_1f(\sigma(g_2,g_3))=g_1\partial f(g_2,g_3)$.
\par
It must be checked that cohomologous 
1-cocycles in $Z^1(E^\times,i^*A)^\Gamma=Hom_\Gamma(E^\times,i^*A)$ 
are mapped to cohomologous 2-cocycles in 
$Z^2(\Gamma,A,\overline{\alpha})$. (We
use that fact, which we have essentially already verified,
that $\Gamma$-invariant classes 
may be represented by $\Gamma$-invariant cocycles.)
A 2-cocycle $(f,{\alpha},\overline{\alpha})$ is
{\bf cohomologous} to a 2-cocycle 
$(f',{\alpha}',\overline{\alpha})$
if there is a map $c:\Gamma\rightarrow A$ for which
$$
f'(g_1,g_2)=
c(g_1g_2)f(g_1,g_2)\alpha(g_1)c(g_2)^{-1}c(g_1)^{-1},
$$
for all $g_j\in\Gamma$, and
$$
\alpha (g)(a)=c(g)\alpha(g)(a)c(g)^{-1},
$$
for all $g\in \Gamma$, $a\in A$. This defines an equivalence 
relation on $Z^2(\Gamma,A,\overline{\alpha})$. (Note that the 
2-cocycles here differ from those of Springer by an inversion.)
The set of
equivalence classes is the second cohomology set of
$\Gamma$ with values in $A$ (with $\Gamma$-group structure
``twisted'' by $\overline{\alpha}$), denoted
$H^2(\Gamma,A,\overline{\alpha})$.
\par
Let $f,f'$ be two cohomlogous 1-cocycles in 
$Z^1(E^\times,i^*A)^\Gamma$. There is an $a\in A$ such that
$f'(x)=a^{-1}f(x)a$, for all $x\in E^\times$.
Since $f,f'$ are $\Gamma$-equivariant,
$$
f'(gx)=gf'(x)=ga^{-1}gf(x)ga=ga^{-1}f(gx)ga,
$$
for all $x\in E^\times$ and all $g\in \Gamma$.
This forces $a\cdot ga^{-1}$ to belong to the centralizer of 
the group $f(E^\times)$ in $A$. (In particlar, such an $a\in A$ 
defines a class $g\longmapsto a\cdot ga^{-1}$ in 
$H^1(\Gamma,Cent_A(f(E^\times)))$, for what that's worth.)
Now we prove that, for cohomologous 1-cocycles $f,f'$ as above,
their images $\partial f,\ \partial f'$ in 
$Z^2(\Gamma,A,\overline{\alpha})$ will be cohomologous.
This is Proposition 1.19 in Springer's article \cite{Spr}. 
Thus, $\partial$ defines a connecting map
from $H^1$ to $H^2$, as desired.
\vskip .1in
{\bf Semi-exactness at $H^1(E^\times,i^*A)$}
\vskip .1in
Let $f\in H^1(W,p^*A)$. We claim that $\partial (res(f))\in 
H^2(\Gamma,A)$ is the trivial 2-cocycle.
\par
If $f'=res(f)$ then 
$$
\split
\partial f'(g_1,g_2)
&=f(s(g_1g_2)s(g_2)^{-1}s(g_1)^{-1})\\
&=f(s(g_1g_2))s(g_1g_2)f(s(g_2)^{-1}s(g_1)^{-1})\\
&=f(s(g_1g_2))w_1w_2f(w_2^{-1}w_1^{-1})\\
&=f(s(g_1g_2))f(w_1w_2)^{-1},
\endsplit
$$
where $g_i=p(w_i)$. Now use the 1-cocycle condition to see that
this is cohomologous to the trivial 2-cocycle in the sense
above, with $c(g)=f(s(g))$.
From this semi-exactness follows and the lemma is completely 
proven.
\vskip .1in
{\bf Conjecture: Exactness at $H^1(E^\times,i^*A)$}
\vskip .1in
Consider an element $f\in Z^1(E^\times,i^*A)^\Gamma$
for which the class of $\partial f$ in $H^2(\Gamma,A,\overline{\alpha})$
is trivial. 
We conjecture that there is an $f'\in H^1(W,p^*A)$ such that 
$f=res(f')$. Since $\partial f$ is trivial, there is a 
$c:\Gamma\rightarrow A$ such that
$$
\partial f(g_1,g_2)=f(s(g_1g_2)s(g_2)^{-1}s(g_1)^{-1})=
c(g_1g_2)g_1c(g_2)^{-1}c(g_1)^{-1},
$$
for all $g_i\in \Gamma$. Let
$$
f'(w):=
\cases
f(x),&               w=i(x)\in i(E^\times),\\
c(g)f(ws(g)^{-1}),&  w\in W,\ p(w)=g,
\endcases 
$$
where $s$ is the section choosen in the statement of the lemma.
By exactness, 
$$
f'(wi(x))=c(g)f(wi(x)s(g)^{-1})=c(g)f(ws(g)^{-1}i(gx))
=f'(w)f(gx)
$$
and 
$$
f'(i(x)w)=c(g)f(i(x)ws(g)^{-1})=f'(w)f(x),
$$ 
where $p(w)=g$, $x\in E^\times$. In particular, $f'(wi(x))=
f'(i(gx)w)$. We conjecture that $f'$ defines a 1-cocycle in
$Z^1(W,p^*A)$.             
\vskip .3in
{\bf References}
\vskip .2in
\item{[Se]} J.-P. Serre, \underbar{Local} \underbar{Fields},
Springer-Verlag Grad. Texts, 1979
\item{[Spr]} T. A. Springer, ``Nonabelian $H^2$ in Galois
cohomology'', in Proc. Symp. Pure Math., vol. 9,
(\underbar{Algebraic} \underbar{Groups}
\underbar{and} \underbar{Discontinuous} \underbar{Subgroups},
Boulder conf., 1965, ed. by A. Borel and G. Mostow), AMS, 1966
\vskip .3in
David Joyner, Dept Math, USNA, Annapolis, MD 21402
\vskip .1in

\enddocument




